\section{The exact bridge at a trivial zero}
\label{bridge:section}

The derivative-tower draft states the nonvanishing logarithmic bridge
correctly, but describes its trivial-zero analogue only as a ratio of
leading nonzero derivatives.  The precise coefficient matters: the
second-to-first derivative ratio appears with a factor of $1/2$.
The following statement supplies the omitted formula and proof.

\begin{theorem}[Regularized harmonic-number bridge]\label{bridge:trivial}
Let $\chi$ be a primitive Dirichlet character of conductor $q$, including
the conductor-$1$ character, and let $k\geq1$.  If
$\chi(-1)=(-1)^k$, then $L(s,\chi)$ has a simple zero at $s=-k$ and
\begin{equation}\label{bridge:formula}
 \frac{L'(k+1,\chi)}{L(k+1,\chi)}
 +\frac12\overline{\left(
       \frac{L''(-k,\chi)}{L'(-k,\chi)}\right)}
 =\gamma+\log\frac{2\pi}{q}-H_k.
\end{equation}
If instead $\chi(-1)=(-1)^{k+1}$, then $L(-k,\chi)\ne0$ and the
corresponding identity is
\begin{equation}\label{bridge:nontrivial}
 \frac{L'(k+1,\chi)}{L(k+1,\chi)}
 +\overline{\left(
       \frac{L'(-k,\chi)}{L(-k,\chi)}\right)}
 =\gamma+\log\frac{2\pi}{q}-H_k.
\end{equation}
\end{theorem}

\begin{proof}
Write $\chi(-1)=(-1)^\delta$, $\delta\in\{0,1\}$, and
\[
 \Lambda(s,\chi)=
 \left(\frac q\pi\right)^{(s+\delta)/2}
 \Gamma\left(\frac{s+\delta}{2}\right)L(s,\chi).
\]
The primitive functional equation \cite{DLMF25} implies, at real regular points,
\begin{equation}\label{bridge:completed}
 \frac{\Lambda'}{\Lambda}(s,\chi)
 =-\overline{\left(
       \frac{\Lambda'}{\Lambda}(1-s,\chi)\right)}.
\end{equation}
The positive value $L(k+1,\chi)$ is nonzero by its absolutely convergent
Euler product.  The functional equation therefore makes
$\Lambda(-k,\chi)$ finite and nonzero.  This also holds for the
conductor-$1$ character: the completed Riemann zeta function is
meromorphic, but neither $-k$ nor $k+1$ is one of its poles.

Suppose first that $\delta\equiv k\pmod2$, and put
$m=(k-\delta)/2\geq0$.  The gamma factor has a simple pole at $-k$,
so $L(s,\chi)$ must have a simple zero there.  With
$a=L'(-k,\chi)\ne0$ and $b=L''(-k,\chi)$, write $s=-k+u$ and expand:
\[
 L(-k+u,\chi)=a u+\tfrac12 b u^2+O(u^3),
\]
\[
 \Gamma(-m+u/2)=\frac{(-1)^m}{m!}
       \left(\frac2u+H_m-\gamma+O(u)\right).
\]
Taking the logarithmic derivative after cancellation gives
\begin{equation}\label{bridge:regularized-log}
 \frac{\Lambda'}{\Lambda}(-k,\chi)
 =\frac12\log\frac q\pi
  +\frac12(H_m-\gamma)+\frac{b}{2a}.
\end{equation}
At the positive point, ordinary differentiation gives
\[
 \frac{\Lambda'}{\Lambda}(k+1,\chi)
 =\frac12\log\frac q\pi
  +\frac12\psi\left(\frac{k+1+\delta}{2}\right)
  +\frac{L'(k+1,\chi)}{L(k+1,\chi)}.
\]
Insert these two formulas in \eqref{bridge:completed}.  The elementary
half-integer digamma identity \cite{DLMF5}
\[
 \psi\left(\frac{k+1+\delta}{2}\right)+H_m-\gamma
 =2H_k-2\gamma-2\log2
\]
follows from $\psi(r+\tfrac12)=-\gamma-2\log2+2H_{2r}-H_r$,
checking $\delta=0$ and $\delta=1$.  It proves
\eqref{bridge:formula}, including the factor $1/2$.

In the opposite parity the gamma factor is finite and nonzero at
$-k$, hence $L(-k,\chi)\ne0$.  Applying
\eqref{bridge:completed} directly gives the sum of the two logarithmic
$L$-derivatives as
\[
 -\log\frac q\pi
 -\frac12\psi\left(\frac{k+1+\delta}{2}\right)
 -\frac12\psi\left(\frac{\delta-k}{2}\right).
\]
The second digamma argument is a half-integer.  Digamma reflection
there has zero cotangent term, and duplication reduces the sum to
$2H_k-2\gamma-2\log2$, proving
\eqref{bridge:nontrivial}.
\end{proof}

For example, the trivial-zero cases $\zeta(-2)=0$ and $\beta(-1)=0$
give respectively
\[
 \frac{\zeta'(3)}{\zeta(3)}
 +\frac{\zeta''(-2)}{2\zeta'(-2)}
 =\gamma+\log(2\pi)-\frac32,
\]
\[
 \frac{\beta'(2)}{\beta(2)}
 +\frac{\beta''(-1)}{2\beta'(-1)}
 =\gamma+\log\frac\pi2-1.
\]
Thus the positive-side first derivative really corresponds to a
negative-side second derivative at the trivial-zero parity.  It cannot
in general be recovered from the negative first-jet value alone.
The accompanying numerical check includes these two identities and
complex primitive quartic characters modulo $5$, for which the
conjugation in \eqref{bridge:formula} is essential.
